\section{Feature denoising}

\subsection{Motivation}

\headingindent
Feature denoising is not classical logit distillation. It follows the
predictive-representation idea associated with joint-embedding predictive
architectures: predict a target representation from a partial or corrupted
context instead of reconstructing raw pixels
\cite{lecun2022path,assran2023ijepa}. Here the teacher remains fixed. A small
student learns an ID-conditional prediction rule in the teacher's feature
space. An OOD input should be harder to reconstruct because its feature
relations were absent from student training.

For a clean teacher feature $z$, a corruption produces $\widetilde z$. The
student predicts $\widehat z=S(\widetilde z)$ and minimizes an MSE on all or
only hidden coordinates. We used three reconstruction scores:
\begin{align*}
s_{\mathrm{rec}}&=-\|\widehat z-z\|_2^2,\\
s_{\mathrm{abs}}&=\|\widetilde z-z\|_2^2-\|\widehat z-z\|_2^2,\\
s_{\mathrm{cos}}&=\frac{\widehat z^\top z}{\|\widehat z\|_2\|z\|_2}.
\end{align*}
These are respectively the negative reconstruction error, the absolute
improvement over the zero-filled or otherwise corrupted input, and cosine
similarity. The results below use absolute improvement because it directly
measures how much the trained denoiser corrects the corruption.

\subsection{Tested variants}

\headingindent
The experiments predicted clean pooled embeddings from affine-augmented,
PixMix, or block-masked images, and reconstructed masked channels, masked
spatial positions, or ViT patch tokens. ResNet feature-map denoisers used
residual CNNs. The ViT patch-token model received a $14\times14$ grid of
768-dimensional layer-12 tokens and used three $3\times3$ convolutions
$768\to64\to64\to768$. A Bernoulli mask hid 20\% of patch tokens, and the loss
was evaluated only at hidden positions.

\begin{table}[ht]
\centering
\rowcolors{2}{rowblue!35}{white}
\begin{tabularx}{\textwidth}{@{}L{3.0cm}L{3.4cm}YY@{}}
\toprule
Variant & Corruption & Student target & Score candidates \\
\midrule
Pixel-to-embedding & Affine or PixMix image & Clean pooled feature &
Negative reconstruction error, absolute improvement, cosine similarity \\
Channel masking & Hidden feature-map channels & Clean feature map &
Negative reconstruction error, absolute improvement \\
Spatial masking & Hidden positions or blocks & Clean feature map &
Negative reconstruction error, absolute improvement \\
ViT patch masking & Hidden layer-12 patch tokens & Clean patch-token grid &
Negative reconstruction error, absolute improvement \\
\bottomrule
\end{tabularx}
\caption{Feature-denoising variants retained in this report.}
\label{tab:feature-denoising-variants}
\end{table}

\subsection{Layerwise results on the initial CIFAR protocol}

\headingindent
For the ResNet-18 channel-masking pipeline, each channel was independently
hidden with probability $0.2$, and a residual CNN reconstructed the clean
feature map. We evaluated the performance on several internal teacher layers: layers 2, 3, and 4. Layer 3
was strongest overall for CIFAR-10 as ID and gave particularly strong SVHN
separation; layer 4 was strongest on MNIST. With CIFAR-100 as ID, layer 3 gave
the best compromise, despite remaining weak on Near-OOD.

\begin{table}[ht]
\centering
\small
\rowcolors{2}{rowblue!35}{white}
\begin{tabularx}{\textwidth}{@{}L{1.75cm}L{1.1cm}*{4}{>{\centering\arraybackslash}X}@{}}
\toprule
ID dataset & Layer & MNIST & SVHN & Near-OOD CIFAR & Macro \\
\midrule
CIFAR-10 & 2 & \metric{0.597}{0.918} & \metric{0.937}{0.271} &
\metric{0.542}{0.913} & \metric{0.655}{0.754} \\
CIFAR-10 & 3 & \metric{0.846}{0.835} & \metric{0.962}{0.217} &
\metric{0.841}{0.662} & \best{\metric{0.872}{0.594}} \\
CIFAR-10 & 4 & \metric{0.979}{0.115} & \metric{0.893}{0.458} &
\metric{0.773}{0.740} & \metric{0.855}{0.513} \\
\midrule
\mbox{CIFAR-100} & 2 & \metric{0.881}{0.545} & \metric{0.900}{0.421} &
\metric{0.552}{0.950} & \metric{0.721}{0.716} \\
\mbox{CIFAR-100} & 3 & \metric{0.994}{0.022} & \metric{0.936}{0.283} &
\metric{0.504}{0.967} & \best{\metric{0.735}{0.560}} \\
\mbox{CIFAR-100} & 4 & \metric{0.621}{0.981} & \metric{0.806}{0.760} &
\metric{0.774}{0.825} & \metric{0.744}{0.848} \\
\bottomrule
\end{tabularx}
\caption{Residual-CNN results with uniform 20\% channel masking and absolute
improvement as the OOD Score. Cells report AUROC / FPR@95 under the original
ID-positive convention.}
\label{tab:channel-denoising-results}
\end{table}

\subsection{Combining layer scores}

\headingindent
Since layer 3 performed better on Far-OOD and layer 4 performed better on Near-OOD, 
it motivated a composition
of their absolute-improvement scores. This experiment was inspired by HEAT,
which composes several energy terms through a generalized energy operator
\cite{lafon2023heat}. We standardized each layer score using its mean and
population standard deviation on ID data, then computed
\[
 s_\beta(z_3,z_4)=\frac{1}{\beta}
 \log\left(e^{\beta z_3}+e^{\beta z_4}\right),\qquad \beta\ne0,
 \qquad
 s_0=\frac{z_3+z_4}{2}.
\]
Negative $\beta$ emphasizes the smaller standardized score, positive $\beta$
emphasizes the larger score, and values near zero produce a smoother balance.

\begin{table}[ht]
\centering
\small
\rowcolors{2}{rowblue!35}{white}
\begin{tabularx}{0.76\textwidth}{@{}L{2.0cm}*{2}{>{\centering\arraybackslash}X}@{}}
\toprule
$\beta$ & CIFAR-10 ID macro & CIFAR-100 ID macro \\
\midrule
$-100$ & \metric{0.893}{0.513} & \metric{0.830}{0.547} \\
$-10$  & \metric{0.895}{0.511} & \metric{0.830}{0.546} \\
$-1$   & \metric{0.916}{0.400} & \metric{0.836}{0.512} \\
$-0.1$ & \metric{0.920}{0.359} & \best{\metric{0.831}{0.502}} \\
$0$    & \metric{0.920}{0.359} & \metric{0.829}{0.505} \\
$0.1$  & \best{\metric{0.920}{0.358}} & \metric{0.827}{0.510} \\
$1$    & \metric{0.917}{0.356} & \metric{0.802}{0.581} \\
$10$   & \metric{0.902}{0.424} & \metric{0.762}{0.697} \\
$100$  & \metric{0.901}{0.440} & \metric{0.760}{0.704} \\
\bottomrule
\end{tabularx}
\caption{HEAT-inspired composition of ResNet-18 layer-3 and layer-4 absolute-
improvement scores. Cells report balanced-macro AUROC / FPR@95 under the
ID-positive convention.}
\label{tab:heat-beta-r18}
\end{table}

The useful range was concentrated near zero. With CIFAR-10 as ID,
$\beta=0.1$ reached \metric{0.920}{0.358}; with CIFAR-100 as ID,
$\beta=-0.1$ reached \metric{0.831}{0.502}, a better performance than on each layer separately without the composition.

\subsection{ViT patch-token denoising}

\headingindent
We tested ViT patch-token denoising quickly during the final days of the
internship, after moving to the OpenOOD-inspired evaluation described in
Section~\ref{sec:openood-transition}. At that point, we also changed the metric
convention from ID-positive to OOD-positive so that the evaluation matched the
OpenOOD benchmark convention.
% These results should therefore be read as a
% late-stage exploratory experiment rather than as a directly comparable part
% of the initial CIFAR study.

Patch-token denoising remained below the teacher-only ViT MSP baseline. The
table reports absolute improvement for the denoiser, consistently with the
ResNet reconstruction results above.

\clearpage
\begin{landscape}

\begin{table}[H]
\centering
\scriptsize
\setlength{\tabcolsep}{2.7pt}
\renewcommand{\arraystretch}{1.12}
\caption{CIFAR-10 as the ID dataset. ViT patch-token denoising and the
teacher-only MSP baseline on OpenOOD v1.5; cells report AUROC / FPR@95 under
the stated positive-class convention.}
\begin{tabularx}{\linewidth}{@{}L{2.75cm}L{1.15cm}L{2.1cm}*{9}{>{\centering\arraybackslash}X}@{}}
\toprule
Variant & Conv. & OOD Score &
\multicolumn{3}{c}{\cellcolor{neargreen}\textbf{Near-OOD}} &
\multicolumn{5}{c}{\cellcolor{farorange}\textbf{Far-OOD}} &
\cellcolor{macroblue}\textbf{Balanced} \\
\cmidrule(lr){4-6}\cmidrule(lr){7-11}\cmidrule(lr){12-12}
& & &
\cellcolor{neargreen}CIFAR-100 & \cellcolor{neargreen}TIN &
\cellcolor{nearaverage}\textbf{Average} &
\cellcolor{farorange}MNIST & \cellcolor{farorange}Places365 &
\cellcolor{farorange}SVHN & \cellcolor{farorange}Textures &
\cellcolor{faraverage}\textbf{Average} & \cellcolor{macroblue}\textbf{Macro} \\
\midrule
\cellcolor{selectedblue!55}\textbf{Teacher ViT-B/16} & \idpos & MSP &
\nearcell{\metric{0.971}{0.124}} & \nearcell{--} & \nearavg{--} &
\farcell{\metric{0.993}{0.024}} & \farcell{--} &
\farcell{\metric{0.991}{0.035}} & \farcell{--} & \faravg{--} &
\balanced{\metric{0.985}{0.061}} \\
\cellcolor{conventionred}\textbf{Teacher ViT-B/16} & \oodpos & MSP &
\nearcell{\metric{0.971}{0.138}} & \nearcell{\metric{0.991}{0.042}} &
\nearavg{\metric{0.981}{0.090}} &
\farcell{\metric{0.995}{0.020}} & \farcell{\metric{0.995}{0.024}} &
\farcell{\metric{0.991}{0.041}} & \farcell{\metric{0.999}{0.0003}} &
\faravg{\metric{0.995}{0.021}} & \balanced{\metric{0.988}{0.056}} \\
\addlinespace
\cellcolor{conventionred}\textbf{Patch-token denoising L12} & \oodpos &
Absolute improvement &
\nearcell{\metric{0.863}{0.673}} & \nearcell{\metric{0.942}{0.386}} &
\nearavg{\metric{0.903}{0.529}} &
\farcell{\metric{0.429}{0.920}} & \farcell{\metric{0.964}{0.215}} &
\farcell{\metric{0.919}{0.430}} & \farcell{\metric{0.996}{0.012}} &
\faravg{\metric{0.827}{0.394}} & \balanced{\metric{0.865}{0.462}} \\
\bottomrule
\end{tabularx}
\label{tab:vit-patch-results-c10}
\end{table}

\vfill

\begin{table}[H]
\centering
\scriptsize
\setlength{\tabcolsep}{2.7pt}
\renewcommand{\arraystretch}{1.12}
\caption{CIFAR-100 as the ID dataset. ViT patch-token denoising and the
teacher-only MSP baseline on OpenOOD v1.5; cells report AUROC / FPR@95 under
the stated positive-class convention.}
\begin{tabularx}{\linewidth}{@{}L{2.75cm}L{1.15cm}L{2.1cm}*{9}{>{\centering\arraybackslash}X}@{}}
\toprule
Variant & Conv. & OOD Score &
\multicolumn{3}{c}{\cellcolor{neargreen}\textbf{Near-OOD}} &
\multicolumn{5}{c}{\cellcolor{farorange}\textbf{Far-OOD}} &
\cellcolor{macroblue}\textbf{Balanced} \\
\cmidrule(lr){4-6}\cmidrule(lr){7-11}\cmidrule(lr){12-12}
& & &
\cellcolor{neargreen}CIFAR-10 & \cellcolor{neargreen}TIN &
\cellcolor{nearaverage}\textbf{Average} &
\cellcolor{farorange}MNIST & \cellcolor{farorange}Places365 &
\cellcolor{farorange}SVHN & \cellcolor{farorange}Textures &
\cellcolor{faraverage}\textbf{Average} & \cellcolor{macroblue}\textbf{Macro} \\
\midrule
\cellcolor{selectedblue!55}\textbf{Teacher ViT-B/16} & \idpos & MSP &
\nearcell{\metric{0.926}{0.375}} & \nearcell{--} & \nearavg{--} &
\farcell{\metric{0.684}{0.845}} & \farcell{--} &
\farcell{\metric{0.930}{0.414}} & \farcell{--} & \faravg{--} &
\balanced{\metric{0.847}{0.545}} \\
\cellcolor{conventionred}\textbf{Teacher ViT-B/16} & \oodpos & MSP &
\nearcell{\metric{0.926}{0.316}} & \nearcell{\metric{0.938}{0.238}} &
\nearavg{\metric{0.932}{0.277}} &
\farcell{\metric{0.689}{0.662}} & \farcell{\metric{0.946}{0.217}} &
\farcell{\metric{0.932}{0.255}} & \farcell{\metric{0.988}{0.076}} &
\faravg{\metric{0.889}{0.303}} & \balanced{\metric{0.911}{0.290}} \\
\addlinespace
\cellcolor{conventionred}\textbf{Patch-token denoising L12} & \oodpos &
Absolute improvement &
\nearcell{\metric{0.629}{0.889}} & \nearcell{\metric{0.772}{0.806}} &
\nearavg{\metric{0.701}{0.847}} &
\farcell{\metric{0.148}{0.979}} & \farcell{\metric{0.822}{0.717}} &
\farcell{\metric{0.644}{0.812}} & \farcell{\metric{0.935}{0.340}} &
\faravg{\metric{0.638}{0.712}} & \balanced{\metric{0.669}{0.780}} \\
\bottomrule
\end{tabularx}
\label{tab:vit-patch-results-c100}
\end{table}

\end{landscape}
\clearpage
